% =========================================================
% Rational Cauchy Sequences
% =========================================================

\section{Rational Cauchy Sequences}
\label{sec:rational-cauchy}

\subsection{Definitions and Theorems}

The rational numbers form an ordered field and they approximate real numbers densely.  At this
stage it is tempting to believe that they should already be adequate for analysis.  The first serious
obstruction appears when one studies limits of rational sequences.

\begin{definition}[Sequence of rational numbers]
\label{def:rational-sequence}
A \emph{sequence of rational numbers} is a function
\[
x:\mathbb{N}\to\mathbb{Q}.
\]
It is usually written $(x_n)_{n\in\mathbb{N}}$ or simply $x_1,x_2,x_3,\dots$.
\end{definition}

\begin{remark}
Thinking of a sequence as a function is useful because it emphasizes that a sequence is not merely a
list but an indexed family of values.  Later this viewpoint makes it natural to add, subtract, and
compare sequences term by term.
\end{remark}

\begin{definition}[Cauchy sequence in $\mathbb{Q}$]
\label{def:cauchy-rational}
A sequence $(x_n)$ of rational numbers is called a \emph{Cauchy sequence} if
\[
\forall \varepsilon>0\ \exists N\in\mathbb{N}\ \forall m,n\ge N,
\qquad
|x_m-x_n|<\varepsilon.
\]
\end{definition}

\begin{remark}
A Cauchy sequence is one whose terms eventually cluster together as tightly as we wish.  Notice that
this definition does not mention a limit at all; it speaks only about the internal behavior of the
sequence.
\end{remark}

\begin{theorem}[Every Convergent Rational Sequence Is Cauchy]
\label{thm:convergent-implies-cauchy-q}
Every convergent sequence in $\mathbb{Q}$ is a Cauchy sequence.
\end{theorem}

\begin{remark*}[Proof]
\hyperref[prf:convergent-implies-cauchy-q]{Go to proof.}
\end{remark*}

\begin{remark}
The converse is the issue.  In a complete space every Cauchy sequence converges.  The rational
numbers are not complete, and so this converse can fail.
\end{remark}

\begin{remark}[Examples]
\begin{enumerate}
  \item The sequence $x_n=1/n$ is Cauchy because it converges to $0$.
  \item The decimal approximations
  \[
  1,\ 1.4,\ 1.41,\ 1.414,\ 1.4142,\dots
  \]
  form a Cauchy sequence in $\mathbb{Q}$.
\end{enumerate}
\end{remark}

\begin{tcolorbox}[colback=thmbox, colframe=thmborder, arc=2pt,
  left=6pt, right=6pt, top=4pt, bottom=4pt,
  title={\small\textbf{Theorem (There Exist Cauchy Sequences in $\mathbb{Q}$ That Do Not Converge to a Rational Number)}},
  fonttitle=\small\bfseries]
\begin{theorem}[There Exist Cauchy Sequences in $\mathbb{Q}$ That Do Not Converge to a Rational Number]
\label{thm:cauchy-not-convergent-q}
There exist Cauchy sequences in $\mathbb{Q}$ that do not converge to a rational number.
\end{theorem}
\end{tcolorbox}

\begin{remark*}[Standard quantified statement]
\[
\exists (x_n)_{n\in\mathbb{N}}\subseteq \mathbb{Q}\;
\Bigl(
\operatorname{CauchySequence}(x_n)
\land
\forall q\in\mathbb{Q}\;\neg\operatorname{ConvergesTo}(x_n,q)
\Bigr).
\]
\end{remark*}

\begin{remark*}[Predicate reading]
\[
\exists (x_n)\;
\bigl(
\operatorname{RationalSequence}(x_n)
\land
\operatorname{CauchySequence}(x_n)
\land
\neg\exists q\in\mathbb{Q}\;\operatorname{ConvergesTo}(x_n,q)
\bigr).
\]
\end{remark*}

\begin{remark*}[Negated quantified statement]
\[
\forall (x_n)_{n\in\mathbb{N}}\subseteq \mathbb{Q}\;
\bigl(
\operatorname{CauchySequence}(x_n)
\Longrightarrow
\exists q\in\mathbb{Q}\;\operatorname{ConvergesTo}(x_n,q)
\bigr).
\]
\end{remark*}

\begin{remark*}[Negation predicate reading]
\[
\forall (x_n)\;
\bigl(
\operatorname{RationalSequence}(x_n)\land \operatorname{CauchySequence}(x_n)
\Longrightarrow
\exists q\in\mathbb{Q}\;\operatorname{ConvergesTo}(x_n,q)
\bigr)
\]
would be the false completeness claim for $\mathbb{Q}$.
\end{remark*}

\begin{remark*}[Failure modes]
The theorem would fail only if every Cauchy sequence of rationals converged to
some rational limit. The decimal approximations to $\sqrt{2}$ show that this is
not so.
\end{remark*}

\begin{remark*}[Failure mode decomposition]
\[
1,\;1.4,\;1.41,\;1.414,\;1.4142,\dots
\]
is Cauchy in $\mathbb{Q}$, but its limit in $\mathbb{R}$ is $\sqrt{2}$, and
\hyperref[thm:sqrt2-irrational]{Theorem~\ref*{thm:sqrt2-irrational}} says that
$\sqrt{2}\notin\mathbb{Q}$.
\end{remark*}

\begin{remark}
The second example above converges to $\sqrt{2}$ in $\mathbb{R}$, but $\sqrt{2}\notin\mathbb{Q}$.  Thus
$\mathbb{Q}$ already contains enough arithmetic to define the sequence, but not enough numbers to
contain its limit.
\end{remark}

\begin{remark*}[Interpretation]
Internal coherence is not enough. A rational sequence can have terms that
settle down relative to one another and still fail to land on a rational
number. This is the structural incompleteness of $\mathbb{Q}$.
\end{remark*}

\begin{remark*}[Dependencies]
\hyperref[def:cauchy-rational]{Cauchy sequence in $\mathbb{Q}$},
the decimal approximation example above, and
\hyperref[thm:sqrt2-irrational]{Theorem~\ref*{thm:sqrt2-irrational}}.
\end{remark*}

\begin{remark*}[Proof idea]
Take the rational decimal truncations of $\sqrt{2}$. Show that they form a
Cauchy sequence in $\mathbb{Q}$, then use irrationality of $\sqrt{2}$ to rule
out convergence to a rational limit.
\end{remark*}

\begin{remark*}[Proof]
\hyperref[prf:cauchy-not-convergent-q]{Go to proof.}
\end{remark*}

\subsection{Consequences}

\begin{remark}[Motivation for the real numbers]
\label{rem:real-motivation}
One standard way to construct the real numbers is to take equivalence classes of rational Cauchy
sequences, where two Cauchy sequences are considered equivalent if their difference tends to $0$.
In this way the missing limits are formally adjoined to the rationals.
\end{remark}
